% 相对于Lebesgue测度的三元分解：nu=(1/2)lambda|[0,1]+(1/4)delta_{1/2}+(1/4)mu_C。绝对连续部分密度h=1/2；点质量1/4位于1/2；Cantor部分质量1/4集中于Cantor集（画出第3步的8个保留区间，每段承担1/32）。横向[0,1]对应8cm。
\begin{tikzpicture}[
  lab/.style={font=\figurefont\fontsize{8}{10}\selectfont},
  note/.style={font=\FigureNoteFont\fontsize{7.5}{9}\selectfont,text=FigureMuted}]
\def\L{8}\def\sy{1.6}
% ac part
\fill[FigureBlueFill] (0,0) rectangle (\L,0.5*\sy);
\draw[FigureBlue,line width=1.1pt] (0,0.5*\sy)--(\L,0.5*\sy);
% point mass
\draw[FigureRed,line width=1.4pt,-{Stealth[length=5pt]}] (0.5*\L,0)--(0.5*\L,1.25*\sy);
\node[note,anchor=west,text=FigureRed] at (0.5*\L+0.1,1.2*\sy) {点质量$\frac14\delta_{1/2}$};
% axis
\draw[->,FigureStroke,line width=.7pt] (-0.1,0)--(\L+0.4,0) node[lab,right] {$x$};
\draw[->,FigureStroke,line width=.7pt] (0,-0.1)--(0,1.5*\sy);
\draw[FigureStroke] (0,0.5*\sy)--(-0.06,0.5*\sy); \node[lab,left] at (-0.04,0.5*\sy) {$\frac12$};
\node[note] at (0.25*\L,0.25*\sy) {$\nu_{\mathrm{ac}}=h\lambda,\ h=\frac12$};
\foreach \x/\t in {0/0,0.5/\frac12,1/1}{\node[lab,below] at (\x*\L,-0.04) {$\t$};}
% Cantor support
\begin{scope}[yshift=-0.85cm]
  \draw[FigureMuted,line width=.4pt] (0,0)--(\L,0);
  \foreach \a in {0,2,6,8,18,20,24,26}{\draw[FigureGreen,line width=3pt] ({\a/27*\L},0)--({(\a+1)/27*\L},0);}
  \node[note,anchor=north] at (0.5*\L,-0.12) {奇异连续部分$\frac14\mu_C$：质量集中于Cantor集（图为第$3$步的$8$个保留区间）};
\end{scope}
\end{tikzpicture}
